% ============================================================================
%  Section 8: Kết luận và hướng phát triển
% ============================================================================
\section{Kết luận và hướng phát triển}
\label{sec:conclusion}

\subsection{Trả lời câu hỏi nghiên cứu}
\label{subsec:conclusion}

\paragraph{RQ1: Các artifact được báo cáo cho thấy điều gì về độ chính xác?}
Trên 234 mẫu test, \model{LayoutXLM} đạt Macro EM 42,96\% và Macro NES 63,89\% ở chế độ present-only, cao hơn Baseline (40,41\% và 55,47\%). Baseline nhỉnh hơn ở trường ngày nhờ các mẫu định dạng có tính quy tắc cao, trong khi LayoutXLM tốt hơn rõ rệt ở trường tổng tiền và địa chỉ. \model{Donut} được suy luận với giới hạn sinh 768 token nhưng cho kết quả rất thấp với Macro EM 0,00\% và Macro NES 0,08\% trên present-only, do mô hình bị suy sụp (model collapse) khi huấn luyện tự hồi quy từ đầu trên tập biên lai nhỏ.

\paragraph{RQ2: Sự khác biệt về latency và khả năng triển khai là gì?}
Hậu xử lý của Baseline từ OCR cache rất nhẹ, trung bình khoảng 1,08 ms/ảnh. Suy luận và hậu xử lý LayoutXLM từ OCR cache mất trung bình 87,11 ms/ảnh — vẫn nằm trong mức chấp nhận được cho triển khai thực tế. Trong khi đó, Donut suy luận trực tiếp từ ảnh (end-to-end) mất trung bình 5.411,17 ms/ảnh; mức này rất cao và khó áp dụng cho xử lý thời gian thực do phải tự hồi quy sinh chuỗi dài 768 token trên GPU.

\paragraph{RQ3: OCR làm giảm hiệu năng LayoutXLM bao nhiêu?}
Thí nghiệm Oracle OCR đã đo được mức hao hụt này. Khi dùng văn bản và hộp bao hoàn hảo (ground-truth), LayoutXLM đạt Macro EM present-only 66,97\% (tăng +24,01 điểm \% so với mức 42,96\% khi dùng VietOCR thực tế). Các trường bị ảnh hưởng nặng nhất là địa chỉ (tăng từ 11,66\% lên 64,57\%) và tên cửa hàng (tăng từ 28,12\% lên 64,29\%), cho thấy chất lượng nhận dạng OCR chính là nút thắt cổ chai lớn nhất của pipeline.

\subsection{Đe dọa tính hợp lệ}
\label{subsec:validity-threats}

\begin{enumerate}
  \item \textbf{Nhãn rỗng:} Khi cả prediction và ground-truth rỗng, all-samples EM/NES được tính đúng. Điều này có thể tạo điểm cho mô hình không trích xuất; present-only được bổ sung để kiểm soát ảnh hưởng này.
  \item \textbf{Khác biệt annotation:} Donut dùng toàn bộ dữ liệu JSON, trong khi LayoutXLM chỉ dùng mẫu có hộp bao. Đây là yếu tố nhiễu khi so sánh kiến trúc.
  \item \textbf{Suy sụp mô hình Donut:} Dù lượt suy luận hiện tại dùng giới hạn sinh 768 token, Donut vẫn bị suy sụp (collapse), sinh lặp từ liên tục — phản ánh độ khó của việc huấn luyện end-to-end trên tập dữ liệu nhỏ.
  \item \textbf{Sai lệch tọa độ cục bộ:} Bốn prediction LayoutXLM cũ dùng kích thước ảnh thô với hộp OCR trên ảnh resize. Sensitivity check cho thấy kết luận tổng thể không bị đảo ngược, nhưng artifact chính vẫn được giữ nguyên để bảo toàn kết quả thực nghiệm.
  \item \textbf{Giới hạn token:} LayoutXLM cắt chuỗi ở 512 token và chưa triển khai sliding window dù cấu hình có \code{doc\_stride}; trường cuối biên lai có thể bị mất.
  \item \textbf{Thiếu log huấn luyện Donut:} Log huấn luyện đầy đủ của LayoutXLM đã được bổ sung trong thư mục \code{outputs/training\_logs/}, tuy nhiên Donut vẫn thiếu log chi tiết do chỉ sử dụng checkpoint đóng băng sẵn có.
  \item \textbf{Phân tích lỗi heuristic:} Taxonomy tự động giúp thống kê nhất quán nhưng không chứng minh nguyên nhân. Các lỗi phức tạp cần rà soát ảnh và token-level output thủ công.
  \item \textbf{Thiếu ablation tiền xử lý:} Dù đã có Oracle OCR để đo ảnh hưởng của OCR, các phép thử ablation về cấu hình tiền xử lý ảnh hoặc augmentation vẫn chưa được làm.
\end{enumerate}

\subsection{Đóng góp và hạn chế}
\label{subsec:limitations}

Đồ án cung cấp pipeline dữ liệu thống nhất, ba phương pháp trích xuất, prediction artifact trên cùng test set, evaluator xử lý rõ giá trị rỗng, manifest khóa bằng checksum và taxonomy lỗi có thể tái lập. Về hạn chế, quy mô dữ liệu nhỏ, annotation không đồng nhất, checkpoint Donut bị model collapse, bốn mẫu LayoutXLM có sai lệch tọa độ, và bằng chứng huấn luyện Donut chưa đầy đủ.

\subsection{Hướng phát triển}
\label{subsec:future_work}

\begin{enumerate}
  \item Tạo sliding windows cho tài liệu dài, gộp prediction giữa các cửa sổ và kiểm tra tác động lên \field{total}/\field{address}.
  \item Khảo sát các kỹ thuật chống suy sụp mô hình cho Donut (như pre-training trên tập CORD lớn hơn, tinh chỉnh nhiệt độ lấy mẫu temperature, áp dụng repetition penalty khi sinh).
  \item Chạy lại ablation tiền xử lý trên danh sách validation cố định, báo cáo độ phủ cache và confidence interval.
  \item Lưu model card cho từng checkpoint gồm commit, môi trường, seed, dữ liệu, thời gian và checkpoint selection criterion.
\end{enumerate}
